% !TeX root = ../../Book1.tex
% 纲要标题：### C.3 群展示、字问题与可判定性
% 纲要位置：计划纲要.md，第 869 行。
% BEGIN APPENDIX OUTLINE SYNC
% #### C.3.1 由关系字推导等式
% #### C.3.2 标准形怎样给出判定算法
% #### C.3.3 不可判定性的含义
% END APPENDIX OUTLINE SYNC

\section{群展示、字问题与可判定性}
\label{b1:sec:C03}

群展示给出对象的定义，但读出一个字是否为单位，是另一个问题。一个定义可以很短，而判断其中的等式是否成立仍然困难。\subsecref{b1:subsec:0804:03}已经给出自由群和二面体群的算法；这里解释这些算法为何有效，以及有限展示本身为什么不保证类似算法存在。

\ProseParagraphBreak
给定群 \(G\) 的一个有限生成集。若存在一个固定算法，对该生成集及其逆组成的每个有限输入字都在有限步后停止，并正确回答它在 \(G\) 中是否表示单位，则称这个生成集下的字问题\term[kepanding]{可判定}[decidable]。停机所用的步数可以随输入增长，也没有规定算法必须很快，但肯定和否定两类输入都必须得到答案。下面先说明：即使不能作出这种总会结束的判断，成立的等式仍然有可以检查的有限推导。

\subsection{由关系字推导等式}
\label{b1:subsec:C03:01}

固定有限展示 \(G=\langle S\mid R\rangle\)，即 \(S\) 是有限生成符号集，\(R\) 是有限个关系字组成的集合。输入是字母表 \(S\sqcup S^{-1}\) 上的有限字 \(w\)，要判断它在 \(G\) 中是否等于单位。判断两个字 \(u,v\) 是否相等，等价于对 \(uv^{-1}\) 作这项判断。\cref{b1:def:group-presentation}说明，\(w=1\) 当且仅当 \(w\) 在自由群 \(F(S)\) 中属于 \(R\) 的正规闭包。

\ProseParagraphBreak
这里须分清两种等式。在自由群中，只准消去相邻的互逆字母；在展示定义的商群中，还可使用关系。例如 \(r=aba^{-1}b^{-1}\) 在自由群 \(F(a,b)\) 中是非空约化字，却在 \(\langle a,b\mid r=1\rangle\) 中等于单位。所谓等式的有限证书，就是把“使用了哪些关系、在哪些字的两侧使用”写成自由群中可以直接检查的等式。

\begin{proposition}\label{b1:prop:C-word-certificate}
设 \(G=\langle S\mid R\rangle\) 为有限展示群，\(w\in F(S)\) 为生成字母及其逆组成的字。字 \(w\) 在 \(G\) 中表示单位，当且仅当存在有限个字 \(u_i\in F(S)\)、关系字 \(r_i\in R\) 及符号 \(\varepsilon_i\in\{1,-1\}\)，使
\begin{equation}\label{b1:eq:C-word-certificate}
w=\prod_{i=1}^m u_i r_i^{\varepsilon_i}u_i^{-1}
\quad\text{在自由群 }F(S)\text{ 中成立}.
\end{equation}
给出这些数据后，可以用有限次自由消去检验等式。
\end{proposition}

\begin{proof}
所有右端乘积（包括空乘积）构成包含 \(R\) 的正规子群：乘法串接因子，取逆颠倒因子顺序并改变符号，用字 \(v\) 共轭只需将各 \(u_i\) 换成 \(vu_i\)。而任意包含 \(R\) 的正规子群都包含这些乘积。因此这个集合恰为正规闭包，给出等价性。检验时把 \(w^{-1}\) 接到右端之前，用\cref{b1:lem:unique-reduced-word}的自由消去算法约化；结果为空字当且仅当自由群中的等式成立。
\end{proof}

这还给出一个只保证找到肯定答案的过程。对 \(N=0,1,2,\ldots\)，枚举 \(m\leq N\)、各 \(u_i\) 长度不超过 \(N\) 的全部证书候选，并逐一检验。每轮候选数有限，任何真实证书又会在足够大的一轮出现。所以若 \(w=1\)，过程终会停机并给出证书；若 \(w\ne1\)，它便一直搜索。这种\term[ban-panding]{半判定}[semidecision]与对每个输入都停机的判定算法不同。把搜索时限用尽，只能说明尚未找到证书，不能证明 \(w\ne1\)。

\ProseParagraphBreak
例如仍取 \(r=aba^{-1}b^{-1}\)，字 \(ara^{-1}\) 在商群中为单位，它的证书只有一个因子，取 \(u_1=a\)、\(r_1=r\)、\(\varepsilon_1=1\) 即可。即使检验者不知道这个证书怎样找到，也能展开并自由消去来确认它。搜索的困难在于事先没有普遍给定的长度界：在某个有限范围内没有找到证书，并不能据此停止全部搜索。

\subsection{标准形怎样给出判定算法}
\label{b1:subsec:C03:02}

若一种标准形不仅存在唯一，而且能在有限步骤内由输入字算出，那么比较标准形就能解决字问题。自由群的逐字消去正具有这些条件；已给乘法表的有限群则可逐字相乘，以单位元作为判据。

\begin{proposition}\label{b1:prop:C-free-product-word-problem}
设有限生成群 \(A,B\) 对各自给定的有限生成集都有可解的字问题。则由两生成集的不相交并生成的 \(A*B\) 也有可解的字问题。
\end{proposition}

\begin{proof}
把输入字分成来自同一因子的最大连续字块。用相应因子的算法判断每个字块是否等于单位；若是，删去它，再把因此相邻的同因子字块串接。每次删除或合并都减少字块数，所以过程终止。停止时，每个剩余字块都表示因子中的非单位元素，且相邻字块来自不同因子。\cref{b1:lem:free-product-normal-form}说明非空的这种交替积不等于单位；空字块列则表示单位。故上述过程对所有输入停机并给出正确答案。
\end{proof}

对于 \(A*_C B\)，\cref{b1:thm:amalgam-normal-form}还使用了陪集代表。如果代表仅由选择公理取得，而没有求分解 \(a=ct\) 的算法，标准形的存在性就不能直接当作字问题的解法。要执行该证明中的化简，至少须能有效完成所用的因子运算、陪集分解及两份 \(C\) 之间的转换，并判断最终 \(C\) 元素是否为单位。

\ProseParagraphBreak
\(G=\langle a,b\mid a^2=b^3\rangle\) 则没有这项困难。记 \(z=a^2=b^3\)，整数除法给出
\[
a^r=z^q a^\epsilon\quad(r=2q+\epsilon,\ 0\leq\epsilon<2),\qquad
b^s=z^p b^\delta\quad(s=3p+\delta,\ 0\leq\delta<3).
\]
\(z\) 与两生成元交换，因此把 \(z\) 的幂移到左端，合并相邻的同因子字块并重复上述除法，即得到 \(z^m t_1\cdots t_n\)。每次合并都缩短剩余字块列；整数运算本身终止。标准形定理说明最终结果等于单位，当且仅当 \(m=0,n=0\)。例如
\[
b^4a^{-3}b^{-1}
=(zb)(z^{-2}a)b^{-1}
=z^{-2}bab^2,
\]
最后一式用了 \(b^{-1}=z^{-1}b^2\)。它是非空标准形，故不等于单位。这里既有存在唯一性，也有求出标准形的具体步骤。

\subsection{不可判定性的含义}
\label{b1:subsec:C03:03}

前面几种算法都利用了各自的特殊结构：自由群有可计算的约化字，有限群有完整乘法表，具体合并积有整数除法给出的陪集分解。是否每个有限展示都至少存在某种字问题算法，即使算法十分缓慢？以下\term[Novikov-Boone-dingli]{Novikov--Boone 定理}[Novikov--Boone theorem]给出否定回答。

\begin{theorem}[Novikov--Boone]\label{b1:thm:C-novikov-boone}
存在一个有限展示的群，其字问题不可判定。
\end{theorem}

\begin{proofsketch}
先取一个固定的、停机集合不可判定的机器 \(T\)。这样的机器可由对角论证得到：若通用机器能判定第 \(n\) 个程序是否在输入 \(n\) 上停机，便可编写一个程序，在判定为停机时无限运行，在判定为不停机时立即停机；把它作用于自己的程序编号，就产生矛盾。

把机器的有限个状态记为 \(q_j\)，带符号记为 \(s_\beta\)。一个带内容和状态组成的字 \(Xq_jY\) 表示一个配置，每条机器指令给出有限个字的替换关系。加上端点符号和停机后的清除关系，得到一个有限展示半群 \(\Gamma\)，使从输入配置到指定终止字 \(q\) 的可达性与停机等价。这里的端点与清除关系须保证逆向使用替换时不会误认停机。

Boone 的群构造再为每条替换添有限个辅助生成元，并加两个生成元 \(k,t\)。令 \(X^\#\) 表示把 \(X\) 中每个字母的指数变号、保留排列顺序所得的字，取特殊字 \(\Sigma=X^\#q_jY\)。构造得到一个固定有限展示群 \(B\)，其中
\[
k(\Sigma^{-1}t\Sigma)k^{-1}(\Sigma^{-1}t^{-1}\Sigma)=1
\quad\Longleftrightarrow\quad Xq_jY=q\text{ 于 }\Gamma.
\]
正向模拟逐条使用群关系，把一次机器步骤化成一次共轭计算；反向论证利用 HNN 标准形或带状关系图，迫使表示单位的上述字出现对应的配置替换，从而恢复一条有限的停机计算。群关系有限，是因为机器状态、带符号和指令都有限。

现在从输入字有效算出上式左端的群字。若 \(B\) 的字问题可判定，就能据此判定 \(T\) 是否停机，矛盾。概要省略半群替换的端点核验及反向关系图的消去证明；它们正是编码中避免伪计算的技术部分，见\cite[Chapter 12, Corollary 12.6; Lemma 12.7; Theorem 12.8, pp. 429--431; proof of Boone's lemma, pp. 432--447]{Rotman1995TheoryGroups}。
\end{proofsketch}

定理中的群和有限展示是固定的，变化的只有输入字；结论因此不只是说“找不到同时处理所有群展示的统一程序”。

\ProseParagraphBreak
有限展示仍保证表示单位的字具有\eqref{b1:eq:C-word-certificate}这样的有限证书。不可判定性否定的是一个对所有输入都能结束的判断过程；它不否定特定字可以被证明非平凡，也不否定自由群、有限群及前述具体合并积的可解算法。继续学习时，可以分别研究哪些群具有有效的标准形、哪些几何条件给出消去算法，以及如何证明一个判断问题本身不可判定。

\ProseParagraphBreak
若还能对所有不表示单位的字设计一个必定确认的搜索过程，两种搜索就能交替各做一步：输入为单位时，肯定搜索最终结束；输入非单位时，否定搜索最终结束。这样就得到一个总会停机的判定算法。因此，对定理中的群，这种否定搜索不可能与肯定证书搜索同时存在。这使“不可判定”与“目前尚未找到一个够快的程序”有了明确区别。
